\centerline{\textbf{\large AGRADECIMIENTOS}}

\vspace{1cm}

% -- Agradecimientos de Alejandro Marín --
En primer lugar, agradezco a Dios, por acompañarme y guiarme a lo largo de este camino, por su
amor incondicional y por brindarme la fortaleza necesaria para superar cada desafío.

A mis padres, Néstor Marín y Gabriela Insfrán, y a mi hermana, Romina Marín, por estar siempre a
mi lado, por confiar en mí desde el primer momento y por brindarme su amor, comprensión y apoyo
incondicional durante toda mi formación. Su esfuerzo, sus palabras de aliento y su acompañamiento
han sido fundamentales para alcanzar esta meta.

A mis abuelos, Ubaldino Marín, Cecilia Moral, Pantaleón Insfrán y Otilia Ortiz, por creer siempre
en mí, incluso en aquellos momentos en los que pocos lo hacían. Gracias por su fe inquebrantable en
mis capacidades, por animarme a seguir adelante y por recordarme constantemente que podía llegar
lejos. Sus palabras de aliento y su confianza han sido una fuente invaluable de motivación a lo
largo de este camino. Asimismo, agradezco a toda mi familia por su cariño, su apoyo y por estar
siempre presente en cada etapa de mi vida.

A la Facultad Politécnica de la Universidad Nacional de Asunción, por abrirme sus puertas y
brindarme un espacio de formación académica y personal, donde pude dar mis primeros pasos en el
mundo de la tecnología y adquirir los conocimientos que hoy constituyen las bases de mi formación
profesional.

A nuestro tutor, PhD. Marcos Villagra, por acompañarnos durante todo el proceso de elaboración de
este trabajo, por su compromiso, su disposición y por la orientación brindada en cada etapa de su
desarrollo.

A mi compañero de tesis, Javier López, por su compromiso, responsabilidad y dedicación constante,
así como por el esfuerzo compartido, el trabajo en equipo y el acompañamiento a lo largo de este
proyecto.

A mis compañeros y amigos de la facultad, Camila Maidana, Rodolfo Vergara, Diego Duarte, Jessica
Alarcón, Guillermo Paiva, Alejandro Adorno, Celia Báez, Guillermo Rodas, Francisco Sanabria, Alejandro Cabral,
Santiago Falcón, Junior Gutiérrez, Matías Franco, Matías Miranda, Arturo Jara, Cristhian Ortellado y
Romina Alfonzo por su amistad, su apoyo y por estar presentes en los momentos en que más los necesitaba.
Gracias por sus palabras de aliento, por motivarme a seguir adelante y por formar parte de una etapa tan importante de mi vida.

Finalmente, a mis compañeros de trabajo y a mis amigos de la iglesia, por su amistad, su cariño y
el apoyo que siempre me brindaron. Gracias por acompañarme en este camino y por hacerme sentir
que, ante cualquier dificultad, siempre podía contar con ustedes.

\vspace{0.5cm}

\rightline{\textit{Alejandro Marín}}

\clearpage

% -- Agradecimientos de Javier López --
\vspace*{1cm}

A mis padres, por ser el pilar fundamental de mi vida, por su amor incondicional, su esfuerzo
silencioso y por darme las herramientas y el respaldo necesarios para alcanzar esta meta. A mis
hermanas, por su compañía constante, su paciencia en los días de mayor exigencia y por celebrar
cada uno de mis logros como propios.

Por sobre todo, agradezco a mi abuela, quien me acompañó desde el cielo en cada momento
difícil que atravesé durante la carrera. En los instantes en que el cansancio parecía ganar, su
recuerdo y sus enseñanzas fueron mi mayor fortaleza. Este título no es más que el principio del
legado de superación y valores que ella inició con sus hijos y continuó con nosotros.

A mi compañero de tesis, Alejandro Marín, por ser un aliado inmejorable en este desafío. Gracias
por todas esas noches de desvelo, por no bajar los brazos ante los obstáculos técnicos y por todos
los cafés que degustamos en el camino mientras depurábamos ideas y código.

A nuestro tutor, el PhD. Marcos Villagra, por su paciencia infinita, su rigor académico y su
vocación docente. Gracias por guiarnos con claridad y por explicarnos más veces de las necesarias
hasta lograr que comprendiéramos a fondo los conceptos abstractos y complejos que conlleva
adentrarse en un área tan desafiante como lo es \textit{blockchain}.

A mis compañeros y amigos de la facultad: Camila Maidana, Mikaela Solís, Mario Gueyraud, Matías
Franco, José Figueredo, Francisco Riquelme, Alejandro Notario, Antonio Insfrán, William González,
Cristhian Ortellado, Guillermo Rodas, Ismael Toledo, José Benítez, Rubén Benítez, Santiago Falcón,
Víctor Ayala, Emiliano Ortiz, Tania Ocampos, Gabriel Villar, Fernando Álvarez y a toda la promo
IIN 019. Gracias por las horas de estudio compartidas, el apoyo mutuo en cada examen y por
demostrarme que la verdadera amistad trasciende las aulas y los pasillos de la facultad para
convertirse en un lazo para toda la vida.

Finalmente, al movimiento «Tu Voz, Nuestra Fuerza», por recibirme con los brazos abiertos desde 
el principio y brindarme un verdadero espacio de pertenencia. Gracias por permitirme crecer en 
el ámbito gremial y social de la facultad. Nunca pierdan esa esencia y sigan demostrando que la 
voz de cada uno es, en realidad, la fuerza que impulsa a todos.

\vspace{0.5cm}

\rightline{\textit{Javier López}}
